% -*- coding: utf-8 -*-
% !TEX program = xelatex
\documentclass[plain,noanswer,medmath]{../../template/jnuexam}
\usepackage{../../template/kysx}

\begin{document}


\examtitle{1991}{4}


\makepart{填空题}{本题满分15分，每小题3分}


\begin{problem}
设$z = \e^{\sin(xy)}$，则$\dz=$ \fillin{}.
\end{problem}


\begin{problem}
设曲线$f(x) = x^3 + ax$与$g(x) = bx^2 + c$都通过点$(-1,0)$，
且在点$(-1,0)$有公共切线，则$a =$ \fillin{}, $b =$\fillin{}, $c =$ \fillin{}.
\end{problem}


\begin{problem}
设$f(x) = x\e^x$，则$f^{(n)}(x)$在点$x=$ \fillin{}处取极小值\fillin{} .
\end{problem}


\begin{problem}
设$A$和$B$为可逆矩阵，$X = \begin{pmatrix}
  0 & A \\
  B & 0
\end{pmatrix}$为分块矩阵，则$X^{-1} =$ \fillin{}.
\end{problem}


\begin{problem}
设随机变量$X$的分布函数为
$$F(x) = P\{X \le x\} = \begin{cases}
    0, & x < - 1, \\
  0.4, & -1 \le x < 1, \\
  0.8, & 1 \le x < 3, \\
    1, & x \ge 3.
\end{cases}$$
则$X$的概率分布为 \fillin{}.
\end{problem}


\makepart{选择题}{本题满分15分，每小题3分}


\begin{problem}
下列各式中正确的是\pickout{}
\begin{abcd}
\item $\lim_{x \to 0^+} \left(1 + \frac{1}{x} \right)^x = 1$．
\item $\lim_{x \to 0^+} \left(1 + \frac{1}{x} \right)^x = \e$．
\item $\lim_{x \to \infty} \left(1 - \frac{1}{x} \right)^x = -\e$．
\item $\lim_{x \to \infty} \left(1 + \frac{1}{x} \right)^{- x} = \e$．
\end{abcd}
\end{problem}


\begin{problem}
设$0 \le a_n < \frac{1}{n}$ ($n = 1,2, \cdots$)，则下列级数中肯定收敛的是\pickout{}
\begin{abcd}
\item $\sum_{n = 1}^{\infty}a_n$．
\item $\sum_{n = 1}^{\infty}(- 1)^n a_n$．
\item $\sum_{n = 1}^{\infty}\sqrt{a_n}$．
\item $\sum_{n = 1}^{\infty}(- 1)^n a_n^2$．
\end{abcd}
\end{problem}


\begin{problem}
设$A$为$n$阶可逆矩阵，$\lambda$是$A$的一个特征根，则$A$的伴随矩阵$A^*$的特征根之一是\pickout{}
\begin{abcd}
\item $\lambda^{-1}|A|^n$．
\item $\lambda^{-1}|A|$．
\item $\lambda |A|$．
\item $\lambda |A|^n$．
\end{abcd}
\end{problem}


\begin{problem}
设$A$和$B$是任意两个概率不为零的不相容事件，则下列结论中肯定正确的是 \pickout{}
\begin{abcd}
\item $\overline{A}$与$\overline{B}$不相容．						
\item $\overline{A}$与$\overline{B}$相容．
\item $P(AB) = P(A)P(B)$．
\item $P(A - B) = P(A)$．
\end{abcd}
\end{problem}


\begin{problem}
对于任意两个随机变量$X$和$Y$，若$E(XY) = E(X) \cdot E(Y)$，则\pickout{}
\begin{abcd}
\item $D(XY) = D(X) \cdot D(Y)$．				
\item $D(X + Y) = D(X) + D(Y)$．
\item $X$和$Y$独立．						
\item $X$和$Y$不独立．
\end{abcd}
\end{problem}


\makepart{}{本题满分5分}

\begin{problem}
求极限$\lim_{x \to 0} \left(\frac{\e^x + \e^{2x} + \cdots + \e^{nx}}{n} \right)^{\frac{1}{x}}$，
其中$n$是给定的自然数．
\end{problem}


\makepart{}{本题满分5分}

\begin{problem}
计算二重积分$I = \iint_D y\dx\dy$，其中$D$是由$x$轴，$y$轴与曲线$\sqrt{\frac{x}{a}} + \sqrt{\frac{y}{b}} = 1$所围成的区域，$a > 0,b > 0$．
\end{problem}


\makepart{}{本题满分5分}

\begin{problem}
求微分方程$xy\frac{\dy}{\dx} = x^2 + y^2$满足条件$\left. y \right|_{x = \e} = 2\e$的特解．
\end{problem}


\makepart{}{本题满分6分}

\begin{problem}
假设曲线$L_1$：$y = 1 - x^2\left(0 \le x \le 1 \right)$，$x$轴和$y$轴所围区域被曲线
$L_2$：$y = ax^2$分为面积相等的两部分，其中$a$是大于零的常数，试确定$a$的值．
\end{problem}


\makepart{}{本题满分8分}

\begin{problem}
某厂家生产的一种产品同时在两个市场销售，售价分别为$p_1$和$p_2$，
销售量分别为$q_1$和$q_2$，需求函数分别为$q_1 = 24 - 0.2p_1$和$q_2 = 10 - 0.05p_2$，
总成本函数为$C = 35 + 40\left(q_1 + q_2 \right)$．
试问：厂家如何确定两个市场的售价，能使其获得的总利润最大？最大利润为多少？
\end{problem}


\makepart{}{本题满分6分}

\begin{problem}
试证明函数$f(x) = \left(1 + \frac{1}{x} \right)^x$在区间$(0, +\infty)$内单调增加．
\end{problem}


\makepart{}{本题满分7分}

\begin{problem}
设有三维列向量
$$\alpha_1 = \begin{pmatrix}
  1 + \lambda \\
  1 \\
  1
\end{pmatrix},\alpha_2 = \begin{pmatrix}
  1 \\
  1 + \lambda \\
  1
\end{pmatrix},\alpha_3 = \begin{pmatrix}
  1 \\
  1 \\
  1 + \lambda
\end{pmatrix},\beta = \begin{pmatrix}
  0 \\
  \lambda \\
  \lambda^2
\end{pmatrix},$$
问$\lambda$取何值时，
\begin{enumerate}
\item $\beta$可由$\alpha_1,\alpha_2,\alpha_3$线性表示，且表达式唯一？
\item $\beta$可由$\alpha_1,\alpha_2,\alpha_3$线性表示，且表达式不唯一？
\item $\beta$不能由$\alpha_1,\alpha_2,\alpha_3$线性表示？
\end{enumerate}
\end{problem}


\makepart{}{本题满分6分}

\begin{problem}
考虑二次型$f = x_1^2+ 4x_2^2+ 4x_3^2+ 2\lambda x_1x_2 - 2x_1x_3 + 4x_2x_3$．
问$\lambda$取何值时，$f$为正定二次型？
\end{problem}


\makepart{}{本题满分6分}

\begin{problem}
试证明$n$维列向量组$\alpha_1,\alpha_2, \cdots ,\alpha_n$线性无关的充分必要条件是
$$D = \begin{vmatrix}
  \alpha_1^T\alpha_1 & \alpha_1^T\alpha_2 & \cdots & \alpha_1^T\alpha_n \\
  \alpha_2^T\alpha_1 & \alpha_2^T\alpha_2 & \cdots & \alpha_2^T\alpha_n \\
              \vdots &             \vdots &        & \vdots \\
  \alpha_n^T\alpha_1 & \alpha_n^T\alpha_2 & \cdots & \alpha_n^T\alpha_n
\end{vmatrix} \ne 0,$$
其中$\alpha_i^T$表示列向量$\alpha_i$的转置，$i = 1,2, \cdots ,n$．
\end{problem}


\makepart{}{本题满分5分}%【类似试卷五第十三题】

\begin{problem}
一汽车沿一街道行驶，需要通过三个均设有红绿信号灯的路口，
每个信号灯为红或绿与其他信号灯为红或绿相互独立，且红绿两种信号显示的时间相等，
以$X$表示该汽车首次遇到红灯前已通过的路口的个数．求$X$的概率分布．
\end{problem}


\makepart{}{本题满分6分}

\begin{problem}
假设随机变量$X$和$Y$在圆域$x^2 + y^2 \le r^2$上服从联合均匀分布．\par
\begin{enumerate*}
\item 求$X$和$Y$的相关系数$\rho$；
\item 问$X$和$Y$是否独立？
\end{enumerate*}
\end{problem}


\makepart{}{本题满分5分}

\begin{problem}
设总体$X$的概率密度为
$$p(x;\lambda) = \begin{cases}
  \lambda ax^{a-1}\e^{-\lambda x^a}, & x > 0, \\
                                  0, & x \le 0,
\end{cases}$$
其中$\lambda > 0$是未知参数，$a > 0$是已知常数．
试根据来自总体$X$的简单随机样本$x_1,x_2, \cdots ,x_n$,
求$\lambda$的最大似然估计量$\hat{\lambda}$．
\end{problem}



\end{document}
